\section{Introducción: por qué este episodio es una clase sobre las rutas tecnológicas de la conducción autónoma}

Esta sección sitúa el episodio. Lang Xianpeng (郎咸朋) trabajó de 2013 a 2018 en Baidu (百度) en conducción autónoma y mapas de alta precisión; en 2018 se incorporó a Li Auto (理想汽车), y durante la última década ha estado en la primera línea de la conducción autónoma en China. El valor de esta entrevista no está en los chismes sobre una empresa, sino en que, desde la perspectiva de alguien que lo vivió, recorre la ruta tecnológica de la conducción autónoma desde «mapas de alta precisión + LiDAR + reglas» hasta «BEV+Transformer», «de extremo a extremo» y «VLM + modelo del mundo + RL».

\begin{importantbox}{Tesis central del episodio}
El hilo conductor de diez años de conducción autónoma es el paso de desarrollar el vehículo como si fuera un «tranvía» a tratar la conducción como una capacidad que se puede aprender. La ruta temprana dependía de mapas de alta precisión, LiDAR costosos e ingeniería de reglas; la contribución clave de Tesla fue resolver el problema elevando la dimensión con visión, BEV, chips y datos de flota; el enfoque de extremo a extremo comprime además las reglas manuales en el modelo y los datos, y empieza a pasar de funciones de software a un sistema de capacidades.
\end{importantbox}

\begin{knowledgebox}{Nota sobre la estrategia visual}
Este video es una entrevista con encuadre fijo, sin proyección de pantalla ni pizarrón. El texto usa la portada solo para identificar la fuente; todas las imágenes del texto son diagramas conceptuales de elaboración propia, que explican la evolución de la pila tecnológica de la conducción autónoma, la elección de sensores, el enfoque de extremo a extremo, los límites de responsabilidad de L3/L4, el ciclo cerrado de datos de Li Auto y la presión organizacional.
\end{knowledgebox}

\begin{figure}[H]
\centering
\includegraphics[width=0.96\textwidth]{av-ten-year-map.png}
\caption{Mapa de diez años de la conducción autónoma: de los mapas de alta precisión y el LiDAR a BEV/Transformer, el enfoque de extremo a extremo y los modelos del mundo. Diagrama conceptual de elaboración propia, basado en la conversación de 00:01:32--00:41:14.}
\end{figure}

\begin{knowledgebox}{Lectura de la figura: de la certeza de ingeniería a la capacidad del modelo}
La figura va de HD Map y LiDAR Stack a BEV, End-to-End y World Model. Al leerla hay que fijarse en el cambio de paradigma: al principio se usaban mapas y sensores para simplificar el mundo; después, datos y modelos para aprender directamente un mundo complejo; más adelante, el sistema necesita predecir el futuro y optimizarse en ciclo cerrado.
\end{knowledgebox}

\subsection{Ruta de lectura}

Esta sección propone cómo leer el episodio. Todo el video puede organizarse en torno a cinco preguntas: primera, ¿por qué la ruta de mapas de alta precisión y LiDAR pesado se parece a un «tranvía»? Segunda, ¿por qué Tesla insiste en la visión y en chips propios? Tercera, ¿qué problema de modelado espacial resolvieron BEV y Transformer? Cuarta, ¿por qué el enfoque de extremo a extremo implica que «antes todos hacíamos mal la conducción autónoma»? Quinta, ¿cómo determinan en conjunto la siguiente etapa L3, L4, los modelos del mundo, RL y el ciclo cerrado de datos y organización de Li Auto?

\noindent\begin{tabularx}{\textwidth}{p{0.19\textwidth}p{0.34\textwidth}X}
\toprule
Pregunta de lectura & Material de la entrevista & Juicio que hay que formarse \\
\midrule
En qué se equivocó la ruta temprana & Mapas de alta precisión, LiDAR, reglas, enumeración exhaustiva de ODD & Trataba las vías abiertas con una mentalidad de rieles y difícilmente alcanza scale. \\
Por qué importa Tesla & Visión pura, BEV, chips, solución por elevación de dimensión & Llevó la conducción autónoma de la acumulación de sensores a una ruta de datos y modelos. \\
Qué cambia el enfoque de extremo a extremo & Los escenarios no pueden enumerarse, las reglas interfieren entre sí, el modelo aprende capacidades & Pasar del desarrollo de funciones al entrenamiento de capacidades. \\
Cómo entender L3/L4 & Responsabilidad del sistema, toma de control, zonas seguras, ODD & L3 no es una extensión de L2; se parece más a un precursor de L4. \\
Cuál es la experiencia de Li Auto & Datos, cómputo, presión organizacional, proyecto Acrópolis (卫城) & La conducción autónoma en producción en serie es un acoplamiento de tecnología, datos, cadena de suministro y organización. \\
\bottomrule
\end{tabularx}

\begin{warningbox}{Límite temporal}
El video se publicó en marzo de 2025, pero la entrevista se realizó en diciembre de 2024. Cuando en la entrevista se dice «este año» suele referirse a 2024, y «el año pasado», a 2023. El texto interpreta según el contexto original de la entrevista y no confunde los tiempos relativos con la fecha de publicación.
\end{warningbox}

\subsection{Resumen de la sección}

EP96 es uno de los episodios más técnicos de la serie de IA e internet de Zhang Xiaojun (张小珺). Puede leerse junto con EP120 sobre Physical AI de XPeng (小鹏), EP132 sobre Galaxea (星海图) y la versión de VLA con proyección de pantalla: la conducción autónoma no es solo una historia de la industria automotriz, sino la avanzada de la IA del mundo físico, los ciclos cerrados de datos, los modelos de extremo a extremo y los modelos del mundo.
